\section{2014 Geometry \& Topology}

\subsection{Individual Exam}

\begin{exercise}
\label{geometry:2014:1}
Let $X$ be the quotient space of $S^{2}$ under the identifications $x \sim -x$ for $x$ in the equator $S^{1}$. Compute the homology groups $H_{n}(X)$. Do the same for $S^{3}$ with antipodal points of the equator $S^{2} \subset S^{3}$ identified.
\end{exercise}

\begin{solution}
\textbf{Prerequisites:} 
\begin{itemize}
    \item \textbf{Quotient Spaces:} Understanding how gluing points together changes the topology of a space.
    \item \textbf{Mayer-Vietoris Sequence:} A tool to compute homology groups by breaking a complex space into two simpler, overlapping spaces.
    \item \textbf{Real Projective Spaces:} Knowing that identifying antipodal points on the boundary of an $n$-disk creates $\mathbb{R}P^n$, and their standard homology groups.
\end{itemize}

\textbf{Intuition:} 
Imagine the 2-sphere $S^2$ as an apple. The equator $S^1$ is the skin around the middle. When we identify antipodal points just on the equator, we are taking opposite points on this middle circle and pulling them together. 

This process transforms the equator into a real projective line, $\mathbb{R}P^1$, which is topologically just a circle. But look at what happens to the top and bottom halves of the apple (the northern and southern hemispheres). Each hemisphere is a disk (or 2-cell) whose boundary is the equator. By identifying opposite points on its boundary, each hemisphere becomes a real projective plane $\mathbb{R}P^2$. 

So, our space $X$ is actually two copies of $\mathbb{R}P^2$ glued together along their common $\mathbb{R}P^1$ boundary. We can use this decomposition to compute the homology using the Mayer-Vietoris sequence!

\textbf{Steps for $X$:}

\textbf{Step 1: Decompose the space.} 
Let $X$ be the space. We can thicken the two halves slightly so they overlap. Let $A$ be the northern hemisphere plus a tiny bit of the southern, and $B$ be the southern hemisphere plus a tiny bit of the northern. Both $A$ and $B$ deformation retract to $\mathbb{R}P^2$. Their intersection $A \cap B$ is a thickened equator with antipodal points identified, which deformation retracts to $\mathbb{R}P^1 \cong S^1$.

\textbf{Step 2: Set up Mayer-Vietoris.}
We apply the Mayer-Vietoris sequence for $X = A \cup B$:
\[ \dots \to H_n(A \cap B) \to H_n(A) \oplus H_n(B) \to H_n(X) \to H_{n-1}(A \cap B) \to \dots \]
We know the homologies:
For $A \cong B \cong \mathbb{R}P^2$: $H_0 = \mathbb{Z}$, $H_1 = \mathbb{Z}_2$, $H_2 = 0$.
For $A \cap B \cong S^1$: $H_0 = \mathbb{Z}$, $H_1 = \mathbb{Z}$, $H_2 = 0$.

\textbf{Step 3: Compute $H_0(X)$.}
Since $X$ is path-connected, $H_0(X) \cong \mathbb{Z}$.

\textbf{Step 4: Compute $H_1(X)$ and $H_2(X)$.}
Look at the sequence ending in degree 1 and 2:
\[ 0 \to H_2(X) \xrightarrow{\partial} H_1(A \cap B) \xrightarrow{(i_*, j_*)} H_1(A) \oplus H_1(B) \to H_1(X) \to 0 \]
Substitute the known groups:
\[ 0 \to H_2(X) \to \mathbb{Z} \xrightarrow{\Phi} \mathbb{Z}_2 \oplus \mathbb{Z}_2 \to H_1(X) \to 0 \]

The map $\Phi: \mathbb{Z} \to \mathbb{Z}_2 \oplus \mathbb{Z}_2$ is induced by the inclusion of $S^1 \to \mathbb{R}P^2$. In $\mathbb{R}P^2$, the generator of the fundamental group (and $H_1$) goes twice around the boundary circle. So the map sends the generator $1 \in \mathbb{Z}$ to $(1, 1) \in \mathbb{Z}_2 \oplus \mathbb{Z}_2$.

The kernel of $\Phi$ is the set of integers mapping to $(0,0)$ mod 2, which are the even integers $2\mathbb{Z} \cong \mathbb{Z}$. Thus, $H_2(X) \cong \ker(\Phi) \cong \mathbb{Z}$.

The image of $\Phi$ is the subgroup generated by $(1,1)$. By the First Isomorphism Theorem and exactness, $H_1(X) \cong (\mathbb{Z}_2 \oplus \mathbb{Z}_2) / \text{Im}(\Phi)$. This quotient has order $4/2 = 2$, so $H_1(X) \cong \mathbb{Z}_2$.

Summary for $X$: $H_0(X) \cong \mathbb{Z}$, $H_1(X) \cong \mathbb{Z}_2$, $H_2(X) \cong \mathbb{Z}$, and $H_n(X) = 0$ for $n > 2$.

\textbf{Steps for $S^3$ quotient $X'$:}

\textbf{Step 5: Decompose $X'$.}
For $S^3$, the "equator" is $S^2$. Identifying antipodal points on this $S^2$ equator turns it into $\mathbb{R}P^2$. 

The space $X'$ consists of two 3-hemispheres (which are 3-balls, $D^3$) glued along their common $\mathbb{R}P^2$ boundary. So $X' = A \cup B$ where $A \cong B \cong \mathbb{R}P^3$ (since an $\mathbb{R}P^3$ is a 3-ball with antipodal boundary points identified). The intersection $A \cap B \cong \mathbb{R}P^2$.

\textbf{Step 6: Apply Mayer-Vietoris for $X'$.}
Homologies:
For $A \cong B \cong \mathbb{R}P^3$: $H_0 = \mathbb{Z}$, $H_1 = \mathbb{Z}_2$, $H_2 = 0$, $H_3 = \mathbb{Z}$.
For $A \cap B \cong \mathbb{R}P^2$: $H_0 = \mathbb{Z}$, $H_1 = \mathbb{Z}_2$, $H_2 = 0$.

Let's look at the sequence for degree 3:
\[ \dots \to H_3(A \cap B) \to H_3(A) \oplus H_3(B) \to H_3(X') \to H_2(A \cap B) \to \dots \]
\[ 0 \to \mathbb{Z} \oplus \mathbb{Z} \to H_3(X') \to 0 \]
This immediately gives $H_3(X') \cong \mathbb{Z} \oplus \mathbb{Z}$.

For lower degrees (degree 1 and 2):
\[ 0 \to H_2(X') \to H_1(A \cap B) \to H_1(A) \oplus H_1(B) \to H_1(X') \to 0 \]
\[ 0 \to H_2(X') \to \mathbb{Z}_2 \xrightarrow{\Psi} \mathbb{Z}_2 \oplus \mathbb{Z}_2 \to H_1(X') \to 0 \]
The map $\Psi: \mathbb{Z}_2 \to \mathbb{Z}_2 \oplus \mathbb{Z}_2$ is induced by inclusion. The $\mathbb{R}P^2$ equator is the standard $\mathbb{R}P^2$ inside $\mathbb{R}P^3$, which induces an isomorphism on $\pi_1$ and $H_1$. So $\Psi$ sends $1 \mapsto (1, 1)$.

The kernel of $\Psi$ is $0$, so $H_2(X') \cong 0$.
The image is the diagonal subgroup, so $H_1(X') \cong (\mathbb{Z}_2 \oplus \mathbb{Z}_2) / \text{Im}(\Psi) \cong \mathbb{Z}_2$.

Summary for $X'$: $H_0(X') \cong \mathbb{Z}$, $H_1(X') \cong \mathbb{Z}_2$, $H_2(X') \cong 0$, $H_3(X') \cong \mathbb{Z}^2$.
\end{solution}

\begin{exercise}
\label{geometry:2014:2}
Let $M \to \mathbb{R}^{3}$ be a graph defined by $z = f(u, v)$ where $\{u, v, z\}$ is a Descartes coordinate system in $\mathbb{R}^{3}$. Suppose that $M$ is a minimal surface. Prove that:
(a) The Gauss curvature $K$ of $M$ can be expressed as
\[ K = \Delta \log\!\left(1 + \frac{1}{W}\right), \qquad W := \sqrt{1 + \left(\frac{\partial f}{\partial u}\right)^{2} + \left(\frac{\partial f}{\partial v}\right)^{2}}, \]
where $\Delta$ denotes the Laplacian with respect to the induced metric on $M$.
(b) If $f$ is defined on the whole $uv$-plane, then $f$ is a linear function (Bernstein theorem).
\end{exercise}

\begin{solution}
\textbf{Prerequisites:} 
\begin{itemize}
    \item \textbf{Minimal Surfaces:} Surfaces that locally minimize area, characterized by having zero mean curvature ($H = 0$).
    \item \textbf{Gauss Map and Curvature:} The relationship between the surface's normal vector, the Gauss map, and the intrinsic Gauss curvature $K$.
    \item \textbf{Laplacian on a Surface:} The Laplace-Beltrami operator $\Delta$ and its properties, particularly for harmonic and superharmonic functions.
\end{itemize}

\textbf{Intuition:} 
Think of a minimal surface as a soap film stretching across a wire frame. Because it minimizes area, the surface is pulled equally in all directions, meaning its principal curvatures are exactly equal and opposite at every point. 

For part (a), the geometry of a surface defined by a graph $z = f(u,v)$ is entirely described by its height function. Because $H=0$, the Gauss curvature $K$ dictates the entire intrinsic geometry. The requested formula elegantly connects this intrinsic curvature to the Laplacian of the normal vector's vertical component.

For part (b), Bernstein's Theorem states a profound global truth: if a soap film is stretched over the entire infinite flat plane, it must itself be a flat plane. A minimal surface has "no room" to wiggle infinitely without eventually turning back on itself (which would violate being a graph over the entire plane). We prove this by showing the surface's curvature forces a certain function to be constant.

\textbf{Steps:}

\textbf{Step 1: Compute the basic geometric quantities.}
The surface $M$ is parameterized by $X(u, v) = (u, v, f(u, v))$.
The tangent vectors are $X_u = (1, 0, f_u)$ and $X_v = (0, 1, f_v)$.
The first fundamental form determinant is $W^2 = EG - F^2 = 1 + f_u^2 + f_v^2$.
The upward-pointing unit normal vector is $N = \frac{1}{W} (-f_u, -f_v, 1)$.
Notice that the vertical (z-axis) component of the normal vector is $N_3 = \frac{1}{W}$.

\textbf{Step 2: Relate the Laplacian of $N_3$ to $K$ (Part a).}
For any minimal surface in $\mathbb{R}^3$, the Gauss map $N: M \to S^2$ satisfies a fundamental eigenvalue equation with respect to the Laplace-Beltrami operator:
\[ \Delta N = -|A|^2 N \]
where $|A|^2$ is the squared norm of the second fundamental form. 
Since the principal curvatures are $k$ and $-k$ (because $H=0$), we have $K = -k^2 \le 0$ and $|A|^2 = k^2 + (-k)^2 = -2K$.
Thus, $\Delta N_3 = 2K N_3$.

Furthermore, the Gauss map is conformal. The norm of the gradient of $N_3$ on $M$ is scaled by $-K$ compared to its gradient on the unit sphere. On the sphere, the gradient squared of the vertical coordinate $z$ is $1 - z^2$. Therefore, on $M$:
\[ |\nabla N_3|^2 = -K(1 - N_3^2) \]

\textbf{Step 3: Evaluate $\Delta \log(1 + N_3)$.}
Let $v = N_3 = 1/W$. We use the chain rule for the Laplacian, $\Delta(g(v)) = g''(v)|\nabla v|^2 + g'(v)\Delta v$:
\[ \Delta \log(1+v) = \nabla \cdot \left( \frac{\nabla v}{1+v} \right) = \frac{\Delta v}{1+v} - \frac{|\nabla v|^2}{(1+v)^2} \]

Substitute our identities for $\Delta v$ and $|\nabla v|^2$:
\[ \Delta \log(1+v) = \frac{2K v}{1+v} - \frac{-K(1 - v^2)}{(1+v)^2} \]
Notice that $1 - v^2 = (1-v)(1+v)$, so the second term simplifies:
\[ \Delta \log(1+v) = \frac{2K v}{1+v} + \frac{K(1-v)}{1+v} = \frac{K(2v + 1 - v)}{1+v} = \frac{K(1+v)}{1+v} = K \]
Since $v = 1/W$, this proves $K = \Delta \log\!\left(1 + \frac{1}{W}\right)$.

\textbf{Step 4: Prove Bernstein's Theorem (Part b).}
Because $M$ is a graph over the entire $uv$-plane, $W = \sqrt{1 + f_u^2 + f_v^2} \ge 1$ everywhere. Thus $v = 1/W$ is strictly bounded: $0 < v \le 1$.
This means the function $\phi = \log(1+v)$ is bounded between $\log(1) = 0$ and $\log(2)$.

Since $K \le 0$ on any minimal surface, we have $\Delta \phi = K \le 0$. This implies $\phi$ is a bounded superharmonic function on $M$.
A deep result in complex analysis and differential geometry states that a global minimal graph over $\mathbb{R}^2$ is conformally equivalent to the complex plane $\mathbb{C}$ (it is a "parabolic" space). 
By Liouville's Theorem, any bounded superharmonic function on $\mathbb{C}$ must be constant.

Therefore, $\phi = \log(1+1/W)$ is a constant function. This forces $W$ to be constant.
If $W$ is constant, its derivatives are zero, meaning $K = \Delta \phi = 0$.
Since $K = 0$ and $H = 0$, both principal curvatures must be zero everywhere. A surface with identically zero principal curvatures is a plane. Thus, $f(u, v)$ must be a linear function.
\end{solution}

\begin{exercise}
\label{geometry:2014:3}
Let $M = \mathbb{R}^{2} / \mathbb{Z}^{2}$ be the two-dimensional torus, $L$ the line $3x = 7y$ in $\mathbb{R}^{2}$, and $S = \pi(L) \subset M$ where $\pi : \mathbb{R}^{2} \to M$ is the projection map. Find a differential form on $M$ which represents the Poincaré dual of $S$.
\end{exercise}

\begin{solution}
\textbf{Prerequisites:} 
\begin{itemize}
    \item \textbf{Torus and Projections:} Understanding $M = \mathbb{R}^2 / \mathbb{Z}^2$ as a square with opposite edges identified, and how lines in $\mathbb{R}^2$ project onto it.
    \item \textbf{Poincaré Duality:} The correspondence between a $k$-dimensional closed curve (cycle) $S$ and an $(n-k)$-dimensional differential form $\eta$ on an $n$-dimensional manifold $M$.
    \item \textbf{Differential Forms and Integration:} Calculating integrals of 1-forms over curves and 2-forms over surfaces.
\end{itemize}

\textbf{Intuition:} 
Imagine water flowing through the torus. The curve $S$ acts like a specific path that winds around the torus 7 times horizontally and 3 times vertically before closing back on itself. 

Poincaré duality asks us to find a "wind" (a differential 1-form $\eta$) that flows everywhere on the torus such that measuring the flux of any other test field $\omega$ against our wind $\eta$ everywhere in the torus gives the exact same result as simply walking along the path $S$ and measuring $\omega$ along the way. In other words, $\eta$ is a smeared-out, continuous version of the discrete curve $S$.

Because the torus is flat and the line is straight, we expect this dual form $\eta$ to be a constant 1-form, $\eta = a\,dx + b\,dy$. We just need to find the correct constants $a$ and $b$ to match the winding numbers.

\textbf{Steps:}

\textbf{Step 1: Parametrize the curve $S$.}
The line $L$ in $\mathbb{R}^2$ is given by $3x = 7y$. We can parameterize it as a path moving from the origin to the next integer point on the line. 
The smallest positive integer solution to $3x = 7y$ is $x = 7, y = 3$.
Thus, a fundamental domain of the line that maps to a single closed loop $S$ on the torus $M$ can be parameterized by:
\[ \gamma(t) = (7t, 3t) \quad \text{for } t \in [0, 1] \]
This means the curve winds 7 times around the $x$-circle and 3 times around the $y$-circle.

\textbf{Step 2: Define the duality condition.}
Let $\eta = a\,dx + b\,dy$ be the Poincaré dual 1-form of $S$. 
The defining property of the Poincaré dual form $\eta$ on a 2-manifold $M$ is that for any closed 1-form $\omega$:
\[ \int_S \omega = \int_M \omega \wedge \eta \]

\textbf{Step 3: Test with the basis form $\omega = dx$.}
First, let's evaluate the line integral over $S$:
\[ \int_S dx = \int_0^1 \frac{dx(\gamma(t))}{dt} dt = \int_0^1 7 \, dt = 7 \]

Next, let's evaluate the surface integral over $M$. We wedge $dx$ with our unknown $\eta$:
\[ \omega \wedge \eta = dx \wedge (a\,dx + b\,dy) = a(dx \wedge dx) + b(dx \wedge dy) = b\,dx \wedge dy \]
Since the total area of the standard torus $M = \mathbb{R}^2 / \mathbb{Z}^2$ is $\int_M dx \wedge dy = 1$:
\[ \int_M b\,dx \wedge dy = b \cdot 1 = b \]

Equating the two results:
\[ b = 7 \]

\textbf{Step 4: Test with the basis form $\omega = dy$.}
Similarly, evaluate the line integral:
\[ \int_S dy = \int_0^1 \frac{dy(\gamma(t))}{dt} dt = \int_0^1 3 \, dt = 3 \]

Now evaluate the surface integral:
\[ \omega \wedge \eta = dy \wedge (a\,dx + b\,dy) = a(dy \wedge dx) + b(dy \wedge dy) = -a(dx \wedge dy) \]
Integrating over the torus:
\[ \int_M -a\,dx \wedge dy = -a \cdot 1 = -a \]

Equating the two results:
\[ -a = 3 \implies a = -3 \]

\textbf{Step 5: Write the final dual form.}
Substitute $a$ and $b$ back into the expression for $\eta$:
\[ \eta = -3\,dx + 7\,dy \]
This 1-form smoothly represents the topological cycle $S$ across the entire torus.
\end{solution}

\begin{exercise}
\label{geometry:2014:4}
Let $p : (\tilde{M}, \tilde{g}) \to (M, g)$ be a Riemannian submersion. This is a submersion $p : \tilde{M} \to M$ such that for each $x \in \tilde{M}$, $dp : \ker^{\perp}(dp) \to T_{p(x)}(M)$ is a linear isometry.
(a) Show that $p$ shortens distances.
(b) If $(\tilde{M}, \tilde{g})$ is complete, so is $(M, g)$.
(c) Show by example that if $(M, g)$ is complete, $(\tilde{M}, \tilde{g})$ may not be complete.
\end{exercise}

\begin{solution}
\textbf{Prerequisites:} 
\begin{itemize}
    \item \textbf{Riemannian Submersions:} Maps between Riemannian manifolds that preserve lengths of horizontal vectors.
    \item \textbf{Horizontal Lifts:} The ability to uniquely lift a curve from the base manifold up to the total space such that its tangent vectors are always horizontal.
    \item \textbf{Hopf-Rinow Theorem:} Understanding metric completeness (Cauchy sequences converge) and geodesic completeness (geodesics extend infinitely).
\end{itemize}

\textbf{Intuition:} 
Imagine looking down at a spiral staircase from directly above. You see a flat circle. The projection from the 3D staircase to the 2D circle is a Riemannian submersion. 

For part (a), if you walk along the staircase, your path involves both horizontal movement (around the circle) and vertical movement (up the stairs). When viewed from above, your shadow only shows the horizontal movement. Because you're ignoring the vertical effort, the shadow's path is shorter than your actual path.

For part (b), if the staircase goes on forever without breaking (is complete), then whenever you trace a geodesic (a straight path) on the circle below, you can uniquely trace a corresponding "horizontal" path on the staircase. Since the staircase is complete, you can walk forever on it, meaning your shadow can walk forever on the circle too. Thus, the circle is complete.

For part (c), completeness of the shadow does not guarantee completeness of the object casting the shadow. A complete circle could be the shadow of a finite, abruptly ending tube.

\textbf{Steps:}

\textbf{Step 1: Prove $p$ shortens distances (Part a).}
Let $\tilde{x}, \tilde{y} \in \tilde{M}$ and let $\tilde{\gamma}(t)$ be any curve connecting them in $\tilde{M}$.
The tangent vector $\tilde{\gamma}'(t)$ can be decomposed into horizontal and vertical components:
\[ \tilde{\gamma}'(t) = \tilde{\gamma}'(t)^H + \tilde{\gamma}'(t)^V \]
where $\tilde{\gamma}'(t)^V \in \ker(dp)$ and $\tilde{\gamma}'(t)^H \in \ker^{\perp}(dp)$.

By the Pythagorean theorem on the inner product $\tilde{g}$:
\[ \|\tilde{\gamma}'(t)\|_{\tilde{g}}^2 = \|\tilde{\gamma}'(t)^H\|_{\tilde{g}}^2 + \|\tilde{\gamma}'(t)^V\|_{\tilde{g}}^2 \ge \|\tilde{\gamma}'(t)^H\|_{\tilde{g}}^2 \]
Since $dp$ is a linear isometry on the horizontal space $\ker^{\perp}(dp)$:
\[ \|dp(\tilde{\gamma}'(t))\|_g = \|dp(\tilde{\gamma}'(t)^H)\|_g = \|\tilde{\gamma}'(t)^H\|_{\tilde{g}} \le \|\tilde{\gamma}'(t)\|_{\tilde{g}} \]
The length of the projected curve $p \circ \tilde{\gamma}$ in $M$ is:
\[ L(p \circ \tilde{\gamma}) = \int \|dp(\tilde{\gamma}'(t))\|_g dt \le \int \|\tilde{\gamma}'(t)\|_{\tilde{g}} dt = L(\tilde{\gamma}) \]
Taking the infimum over all curves connecting $\tilde{x}$ and $\tilde{y}$, we get the distance inequality:
\[ d_M(p(\tilde{x}), p(\tilde{y})) \le d_{\tilde{M}}(\tilde{x}, \tilde{y}) \]

\textbf{Step 2: Prove base manifold completeness (Part b).}
Assume $(\tilde{M}, \tilde{g})$ is complete. By the Hopf-Rinow theorem, it is geodesically complete, meaning any geodesic can be extended for $t \in (-\infty, \infty)$.
Let $\gamma(t)$ be a geodesic in $M$ starting at $x \in M$ with initial velocity $v \in T_x M$.

Choose any point $\tilde{x} \in p^{-1}(x)$. We can find a unique horizontal lift $\tilde{v} \in \ker^{\perp}(dp_{\tilde{x}})$ such that $dp(\tilde{v}) = v$.
Because $\tilde{M}$ is complete, there exists a unique geodesic $\tilde{\gamma}(t)$ in $\tilde{M}$ starting at $\tilde{x}$ with velocity $\tilde{v}$, defined for all $t \in \mathbb{R}$.

A fundamental property of Riemannian submersions is that if a geodesic in $\tilde{M}$ starts with a horizontal tangent vector, its tangent vector remains horizontal for all time, and its projection is a geodesic in $M$.
Therefore, $p \circ \tilde{\gamma}(t)$ is a geodesic in $M$ with initial position $x$ and velocity $v$, and it is defined for all $t \in \mathbb{R}$. Since every geodesic in $M$ can be extended indefinitely, $(M, g)$ is complete.

\textbf{Step 3: Provide a counterexample for the converse (Part c).}
We need a complete base manifold $M$ and an incomplete total space $\tilde{M}$.
Let $M = S^1$, the unit circle, which is compact and therefore complete.
Let $\tilde{M} = (0, 1) \times S^1$, an open cylinder equipped with the standard product metric $g_{Euc} \oplus g_{S^1}$. This space is clearly incomplete (e.g., the sequence $(1/n, \theta)$ has no limit in $\tilde{M}$).

Define $p: (0, 1) \times S^1 \to S^1$ by $p(h, \theta) = \theta$.
This map is a Riemannian submersion because the horizontal space is the $S^1$ direction, on which $dp$ is obviously an isometry.
Thus, $(M, g)$ is complete, but $(\tilde{M}, \tilde{g})$ is not.
\end{solution}

\begin{exercise}
\label{geometry:2014:5}
Let $\Psi : M \to \mathbb{R}^{3}$ be an isometric immersion of a compact surface $M$ into $\mathbb{R}^{3}$. Prove that
\[ \int_{M} H^{2}\,d\sigma \geq 4\pi, \]
where $H$ is the mean curvature of $M$ and $d\sigma$ is the area element of $M$.
\end{exercise}

\begin{solution}
\textbf{Geometric Intuition:} This is the Willmore inequality for the sphere (and a general bound for surfaces). It states that the "average" bending of a closed surface must be at least that of a sphere. The $4\pi$ comes from the fact that the Gauss map of a closed surface must cover the unit sphere at least once, and for a sphere, $H^2 = K$.

\textbf{Proof:} Let $k_1,k_2$ be the principal curvatures and use the convention $H=(k_1+k_2)/2$. Then
\[
H^2-K=\frac{(k_1-k_2)^2}{4}\geq 0.
\]
In particular, on the set $M_+=\{K>0\}$ we have $H^2\geq K$, while outside $M_+$ we still have $H^2\geq 0$. Hence
\[
\int_M H^2\,d\sigma\geq \int_{M_+}K\,d\sigma.
\]

Let $N:M\to S^2$ be the Gauss map. Its Jacobian is $K$ with sign, so the area formula gives
\[
\int_{M_+}K\,d\sigma
\]
as the area of the positively oriented part of the Gauss image counted with multiplicity. Because $M$ is compact, for every $u\in S^2$ a supporting plane with outward normal $u$ touches $\Psi(M)$; at such a support point the Gauss map equals $u$ and $K\geq 0$ in the weak support sense. After a standard small perturbation, this implies that the positive Gauss image covers $S^2$ at least once. Therefore
\[
\int_{M_+}K\,d\sigma\geq \operatorname{area}(S^2)=4\pi.
\]
Combining the two inequalities gives
\[
\int_M H^2\,d\sigma\geq 4\pi.
\]
\end{solution}

\begin{exercise}
\label{geometry:2014:6}
The unit tangent bundle of $S^{2}$ is the subset
\[ T^{1}(S^{2}) = \left\{(p, v) \in \mathbb{R}^{3} \times \mathbb{R}^{3} \mid \|p\| = 1,\; \langle p, v \rangle = 0,\; \|v\| = 1\right\}. \]
Show that it is a smooth submanifold of the tangent bundle $T(S^{2})$ of $S^{2}$ and $T^{1}(S^{2})$ is diffeomorphic to $\mathbb{R}P^{3}$.
\end{exercise}

\begin{solution}
\textbf{Geometric Intuition:} A point in the unit tangent bundle of $S^2$ is an orthonormal pair of vectors $(p, v)$. Adding a third vector $p \times v$ gives a full rotation matrix in $SO(3)$. Thus $T^1(S^2)$ is naturally identified with the rotation group $SO(3)$. Since rotations can be represented by quaternions up to a sign, $SO(3) \cong \mathbb{R}P^3$.

\textbf{Proof:}
Inside $T(S^2)$, a tangent vector is a pair $(p,v)$ with
\[
|p|=1,\qquad \langle p,v\rangle=0.
\]
The unit tangent bundle is the level set
\[
T^1(S^2)=\{(p,v)\in T(S^2): |v|^2=1\}.
\]
The function $(p,v)\mapsto |v|^2$ has differential $2\langle v,\cdot\rangle$ in the fiber direction, which is nonzero on $|v|=1$. Hence $T^1(S^2)$ is a smooth hypersurface of $T(S^2)$.

Define
\[
\Phi:T^1(S^2)\to SO(3),\qquad
\Phi(p,v)=\begin{bmatrix}p&v&p\times v\end{bmatrix}.
\]
The columns form a positively oriented orthonormal basis of $\mathbb{R}^3$, so $\Phi(p,v)\in SO(3)$. Conversely, if $A\in SO(3)$ has columns $(a_1,a_2,a_3)$, then $(a_1,a_2)\in T^1(S^2)$ and
\[
\Phi(a_1,a_2)=A.
\]
Thus $\Phi$ is bijective with smooth inverse $A\mapsto (a_1,a_2)$. Therefore
\[
T^1(S^2)\cong SO(3).
\]

Finally, the standard double cover
\[
S^3\to SO(3)
\]
comes from unit quaternions acting on imaginary quaternions by conjugation. Its kernel is $\{\pm 1\}$, so
\[
SO(3)\cong S^3/\{\pm 1\}\cong \mathbb{R}P^3.
\]
Therefore $T^1(S^2)$ is diffeomorphic to $\mathbb{R}P^3$.
\end{solution}

\subsection{Team Exam}

\begin{exercise}
\label{geometry:2014:7}
Compute the fundamental and homology groups of the wedge sum of a circle $S^{1}$ and a torus $T = S^{1} \times S^{1}$.
\end{exercise}

\begin{solution}
\textbf{Geometric Intuition:} A wedge sum $X \vee Y$ is two spaces joined at a single point. Its fundamental group is the free product (you can loop around one or the other in any order), and its homology is the direct sum (the cycles in each stay independent).

\textbf{Proof:} Let
\[
X=S^1\vee T^2.
\]
The intersection point has a contractible neighborhood in each summand, so Seifert--van Kampen gives
\[
\pi_1(X)\cong \pi_1(S^1)*\pi_1(T^2)
\cong \mathbb{Z}*(\mathbb{Z}\oplus \mathbb{Z}).
\]

For homology, reduced homology of a wedge of connected CW complexes splits:
\[
\widetilde H_k(S^1\vee T^2)
\cong
\widetilde H_k(S^1)\oplus \widetilde H_k(T^2).
\]
Thus
\[
H_0(X)\cong \mathbb{Z},\qquad
H_1(X)\cong \mathbb{Z}\oplus \mathbb{Z}^2\cong \mathbb{Z}^3,
\]
\[
H_2(X)\cong 0\oplus \mathbb{Z}\cong \mathbb{Z},
\qquad
H_k(X)=0\quad (k\geq 3).
\]
\end{solution}

\begin{exercise}
\label{geometry:2014:8}
Given a properly discontinuous action $F : G \times M \to M$ on a smooth manifold $M$, show that $M/G$ is orientable if and only if $M$ is orientable and $F(g, \cdot)$ preserves the orientation of $M$. Use this statement to show that the Möbius band is not orientable and that $\mathbb{R}P^{n}$ is orientable if and only if $n$ is odd.
\end{exercise}

\begin{solution}
\textbf{Geometric Intuition:} Orientation is about a consistent choice of "handedness." If you have a global choice on $M$, it descends to the quotient only if every symmetry in $G$ respects that choice. If a symmetry "flips" the handedness, the quotient becomes non-orientable.

\textbf{Proof of the criterion:}
Let $\pi:M\to M/G$ be the quotient map. Since the action is properly discontinuous, $\pi$ is a covering map locally modeled on diffeomorphisms from open subsets of $M$ to open subsets of $M/G$.

If $M/G$ is orientable, then pulling back its orientation by $\pi$ orients $M$. Every deck transformation $F(g,\cdot)$ preserves this pulled-back orientation because
\[
\pi\circ F(g,\cdot)=\pi.
\]
Thus $M$ is orientable and every group element acts orientation-preservingly.

Conversely, suppose $M$ is orientable and each $F(g,\cdot)$ preserves the chosen orientation. For a point $[p]\in M/G$, choose a small evenly covered neighborhood $U$ and one sheet $\widetilde U\subset M$ mapping diffeomorphically to $U$. Transport the orientation of $\widetilde U$ to $U$. If another sheet is chosen, it differs by some $g\in G$, and the orientation is unchanged because $F(g,\cdot)$ preserves orientation. Hence these local orientations on $M/G$ are well-defined and compatible.

\textbf{Application:}
\begin{enumerate}
    \item The Möbius band is the quotient of $\mathbb{R}\times (-1,1)$ by the action generated by $(x,y)\mapsto (x+1,-y)$. This map has derivative determinant $-1$, so it reverses orientation. Therefore the quotient is not orientable.
    \item $\mathbb{R}P^n=S^n/\{x\sim -x\}$. The antipodal map on $S^n$ has degree $(-1)^{n+1}$, so it preserves orientation exactly when $n$ is odd. Hence $\mathbb{R}P^n$ is orientable if and only if $n$ is odd.
\end{enumerate}
\end{solution}

\begin{exercise}
\label{geometry:2014:9}
(a) Consider the space $Y$ obtained from $S^{2} \times [0, 1]$ by identifying $(x, 0)$ with $(-x, 0)$ and also identifying $(x, 1) \in S^{2} \times \{1\}$ with $(-x, 1)$, for all $x \in S^{2}$. Show that $Y$ is homeomorphic to the connected sum $\mathbb{R}P^{3} \# \mathbb{R}P^{3}$.
(b) Show that $S^{2} \times S^{1}$ is a double cover of the connected sum $\mathbb{R}P^{3} \# \mathbb{R}P^{3}$.
\end{exercise}

\begin{solution}
\textbf{Geometric Intuition:} $\mathbb{R}P^3$ is a 3-ball with antipodal boundary points identified. Taking the connected sum $\mathbb{R}P^3 \# \mathbb{R}P^3$ is like taking a cylinder $S^2 \times [0, 1]$ and "capping" both ends with projective planes.

\textbf{Proof of (a):}
View $\mathbb{R}P^3$ as a closed $3$-ball $D^3$ with antipodal points on $\partial D^3=S^2$ identified. Remove a small open ball around the image of the center of $D^3$. The remaining manifold is homeomorphic to
\[
S^2\times [0,1]
\]
with antipodal identification on one boundary component only. Therefore, taking two copies of this punctured $\mathbb{R}P^3$ and gluing their new spherical boundary components gives exactly
\[
S^2\times[0,1]/\bigl((x,0)\sim(-x,0),\ (x,1)\sim(-x,1)\bigr).
\]
This is the space $Y$, so
\[
Y\cong \mathbb{R}P^3\#\mathbb{R}P^3.
\]

\textbf{Proof of (b):}
Write $S^1$ as $[0,1]/(0\sim 1)$ and define an involution on $S^2\times S^1$ by
\[
\iota(x,[t])=(-x,[-t]).
\]
It is free, since a fixed point would require $x=-x$ on $S^2$, impossible. The quotient has a fundamental domain $S^2\times[0,1/2]$. At the two boundary components the identifications induced by $\iota$ are antipodal on the $S^2$ factor. Hence
\[
(S^2\times S^1)/\langle\iota\rangle
\cong
Y
\cong
\mathbb{R}P^3\#\mathbb{R}P^3.
\]
Thus the quotient map
\[
S^2\times S^1\to \mathbb{R}P^3\#\mathbb{R}P^3
\]
is a double cover.
\end{solution}

\begin{exercise}
\label{geometry:2014:10}
Prove that a bi-invariant metric on a Lie group $G$ has nonnegative sectional curvature.
\end{exercise}

\begin{solution}
\textbf{Geometric Intuition:} Bi-invariant metrics on Lie groups are highly symmetric. The curvature is essentially determined by the Lie bracket $[X, Y]$, which measures the non-commutativity. The positive sectional curvature reflects the fact that these groups are "curved towards themselves" by the group operation.

\textbf{Proof:}
It suffices to compute on left-invariant vector fields. Bi-invariance of the metric implies ad-invariance:
\[
\langle [X,Y],Z\rangle+\langle Y,[X,Z]\rangle=0.
\]
The Koszul formula then gives
\[
\nabla_XY=\frac{1}{2}[X,Y].
\]
Hence
\[
R(X,Y)Y
=\nabla_X\nabla_YY-\nabla_Y\nabla_XY-\nabla_{[X,Y]}Y
=-\frac{1}{4}[Y,[X,Y]].
\]
Using ad-invariance once more,
\[
\langle R(X,Y)Y,X\rangle
=\frac{1}{4}\langle [X,Y],[X,Y]\rangle.
\]
For a two-plane spanned by linearly independent $X,Y$,
\[
K(X,Y)
=
\frac{\langle R(X,Y)Y,X\rangle}
{|X|^2|Y|^2-\langle X,Y\rangle^2}
=
\frac{\|[X,Y]\|^2}
{4\bigl(|X|^2|Y|^2-\langle X,Y\rangle^2\bigr)}
\geq 0.
\]
\end{solution}

\begin{exercise}
\label{geometry:2014:11}
Let $M$ be the upper half-plane $\mathbb{R}_{+}^{2}$ with the metric
\[ ds^{2} = \frac{dx^{2} + dy^{2}}{y^{k}}. \]
For which values of $k$ is $M$ complete?
\end{exercise}

\begin{solution}
\textbf{Geometric Intuition:} Completeness means you can never reach the "edge" in finite time. In the upper half plane, the "edges" are $y=0$ and $y=\infty$. The metric $y^{-k}$ acts as a "speed limit." If $k$ is too large, you reach infinity too fast; if $k$ is too small, you reach zero too fast. $k=2$ is the perfect balance (Hyperbolic space).

\textbf{Proof:}
The metric is conformal to the Euclidean metric:
\[
ds=y^{-k/2}\sqrt{dx^2+dy^2}.
\]
Completeness can fail only by escaping either to the boundary $y=0$ or to infinity.

First consider vertical curves approaching $y=0$:
\[
\gamma(t)=(0,t),\qquad 0<t\leq 1.
\]
Their length near $0$ is
\[
\int_0^1 t^{-k/2}\,dt.
\]
This integral diverges exactly when
\[
\frac{k}{2}\geq 1,
\]
that is, when $k\geq 2$. If $k<2$, the boundary $y=0$ is at finite distance, so the metric is incomplete.

Next consider vertical curves escaping to $y=\infty$:
\[
\gamma(t)=(0,t),\qquad t\geq 1.
\]
Their length is
\[
\int_1^\infty t^{-k/2}\,dt.
\]
This integral diverges exactly when
\[
\frac{k}{2}\leq 1,
\]
that is, when $k\leq 2$. If $k>2$, infinity is at finite distance, so the metric is incomplete.

Thus both ends are infinitely far away only when
\[
k=2.
\]
For $k=2$, the metric is the standard hyperbolic half-plane metric, which is complete. Therefore $M$ is complete exactly for $k=2$.
\end{solution}
